\plainsection{表 1}

我们给每个基本权关联数 1（分别在相应的单表示为正交型时），以及 -1、0（分别在相应表示为辛型、既非正交型也非辛型时）。这个数的计算实质上已在§13中给出，适用于 \(\mathrm{A}_l, \mathrm{B}_l, \mathrm{C}_l, \mathrm{D}_l\) 型。对于 \(\mathrm{E}_6, \mathrm{E}_7, \mathrm{E}_8, \mathrm{F}_4, \mathrm{G}_2\) 型，结果也列于下方（只需应用§7，命题12，以及第 VI 章，§4，第9.VI、9.XI、10.VI、10.XI、11.VI、11.XI、12.VI、12.XI、13.VI、13.XI节）。

\[
\begin{array}{c c} \mathrm{A} _ {l} (l \geq 1) & \mathrm{B} _ {l} (l \geq 2) \\ \varpi_ {r} \left\{ \begin{array}{l l} 0 & r \neq \frac {l + 1}{2} \\ (- 1) ^ {r} & r = \frac {l + 1}{2} \end{array} \right. & \begin{array}{c c} \varpi_ {r} & 1 \quad r \neq l \\ \varpi_ {l} & (- 1) ^ {l (l + 1) / 2} \end{array} \end{array}
\]

\[
\begin{array}{c c} \varpi_ {r} & 1 \quad r \neq l - 1, l \\ \varpi_ {l} \text {and} \varpi_ {l - 1} & \left\{ \begin{array}{l l} 0 & \text {if l is odd} \\ (- 1) ^ {l / 2} & \text {if l is even} \end{array} \right. \end{array}
\]

\[
\begin{array}{c c c c c c c c c c} \mathrm{E} _ {6} & & \mathrm{E} _ {7} & & \mathrm{E} _ {8} & & \mathrm{F} _ {4} & & \mathrm{G} _ {2} \\ \varpi_ {1} & 0 & \varpi_ {1} & 1 & \varpi_ {1} & 1 & \varpi_ {1} & 1 & \varpi_ {1} & 1 \\ \varpi_ {2} & 1 & \varpi_ {2} & - 1 & \varpi_ {2} & 1 & \varpi_ {2} & 1 & \varpi_ {2} & 1 \\ \varpi_ {3} & 0 & \varpi_ {3} & 1 & \varpi_ {3} & 1 & \varpi_ {3} & 1 \\ \varpi_ {4} & 1 & \varpi_ {4} & 1 & \varpi_ {4} & 1 & \varpi_ {4} & 1 \\ \varpi_ {5} & 0 & \varpi_ {5} & - 1 & \varpi_ {5} & 1 \\ \varpi_ {6} & 0 & \varpi_ {6} & 1 & \varpi_ {6} & 1 \\ & & \varpi_ {7} & - 1 & \varpi_ {7} & 1 \\ & & & & \varpi_ {8} & 1 \end{array}
\]
% semantic-fragments: legacy-input-seam
\relax{}
